\chapter{Adenosine Deaminase (ADA)}

\section*{What Is This Test?}
Adenosine deaminase is an enzyme involved in lymphocyte function. ADA activity can be measured in body fluids, particularly pleural fluid, to support evaluation of suspected tuberculosis and selected other infections or inflammatory conditions.

\section*{Alternative Names}
\begin{itemize}
  \item ADA test
  \item Adenosine deaminase activity
  \item Pleural-fluid ADA
  \item Cerebrospinal-fluid ADA
\end{itemize}

\section*{Why This Test Is Ordered}
Pleural-fluid ADA is most useful when tuberculous pleuritis is a consideration, especially in a high-prevalence setting. ADA may also be requested in selected peritoneal, pericardial, or cerebrospinal-fluid evaluations, but it is not a stand-alone diagnostic test.

\section*{How This Test Is Performed}
The laboratory measures ADA activity in a submitted body-fluid specimen using a validated enzymatic method. The specimen type, cell count, protein, glucose, microbiology, and clinical context must accompany the result.

\section*{Preparation, Risks, and Limitations}
The laboratory test has no direct risk; risks depend on how the fluid was obtained. High ADA can occur with bacterial infection, lymphoma, empyema, and other inflammatory conditions. Low ADA does not exclude tuberculosis, particularly in immunosuppression or severe illness.

\section*{Understanding Results}
An elevated pleural-fluid ADA supports tuberculous pleuritis when the clinical and fluid findings are compatible, but culture, molecular testing, histology, and specialist assessment may still be required. Cutoffs vary by specimen type and laboratory.

\section*{Specimen and Preparation}
Pleural-fluid ADA is the most common clinical application. ADA may also be measured in peritoneal (ascitic), pericardial, or cerebrospinal fluid when extrapulmonary tuberculosis is suspected. No fasting is required for the assay, but the fluid must be collected into the correct container, transported promptly, and accompanied by the specimen type and clinical question. Cell count with differential, protein, glucose, lactate dehydrogenase, microbiology, and cytology are often requested on separate portions of the same fluid.

\section*{How Results Are Reported}
Results are commonly reported as activity in U/L. A pleural-fluid value above approximately 40 U/L is often used as a supportive threshold, but the appropriate decision limit depends on local tuberculosis prevalence, the assay, age, and the fluid's cell profile. Values around 40--70 U/L may be indeterminate, whereas markedly higher values are more supportive when the fluid is lymphocyte-predominant and the clinical presentation is compatible. These are interpretive examples, not universal reference intervals.

\section*{False Positives and False Negatives}
ADA can be increased by empyema or complicated bacterial parapneumonic effusion, lymphoma, some other malignancies, rheumatoid pleuritis, and selected inflammatory conditions. Lower concentrations can occur in older adults, immunosuppression, early disease, or severe illness. A low ADA reduces but does not eliminate the possibility of tuberculosis, and a high ADA does not establish it. The result must not be used alone to start or exclude treatment.

\section*{Follow-up and Related Tests}
Evaluation of suspected pleural tuberculosis generally includes acid-fast smear where appropriate, mycobacterial culture, a validated nucleic-acid test when indicated, and clinical assessment. A negative pleural-fluid molecular test or culture does not always exclude pleural disease because organism burden may be low. Depending on the presentation, additional fluid studies, blood tests, or tissue-based testing may be needed. Urgent assessment is required for severe breathlessness, chest pain, confusion, or rapidly worsening illness.

\section*{References}
Centers for Disease Control and Prevention/American Thoracic Society/Infectious Diseases Society of America. Clinical practice guidance for diagnosis of tuberculosis.

World Health Organization. Consolidated guidance on tuberculosis diagnosis.

\clearpage
\section*{Understanding Results and Follow-up}
The laboratory report should identify the specimen, method, units, reference interval or decision limit, and whether the result is positive, negative, abnormal, or inconclusive. Reference intervals vary with age, sex, pregnancy, specimen type, assay, and laboratory; a result outside a range is not automatically a diagnosis, and a result within a range does not always exclude disease. Discuss unexpected results with the ordering clinician, who can decide whether repeat or confirmatory testing, treatment, or referral is appropriate. Do not change medicines or delay urgent care based on an isolated result.


\section*{Regulatory Status and Authorization Date}
This chapter describes a test category, not a specific commercial product. FDA, CDSCO, EU IVDR, and MHRA approval or registration dates therefore cannot be assigned without the assay manufacturer, product name, instrument, and intended use. For laboratory-developed tests, an individual agency approval date may not exist. See the Preface for the interpretation of regulatory status.

\clearpage
